\documentclass[10pt,a4paper]{amsart}
\usepackage{amsmath,amssymb,amsthm,graphicx}
% Retained preamble for the Stacks Project chapter {CH} (shared project preamble).
% The chapter file pinned-original-source.tex only carries \input{{preamble}} on its first line,
% so no preamble text lives inside the chapter itself.  The Stacks Project uses one single shared
% preamble for every chapter: the snapshot stages it three times (sources/stacks-chow/preamble.tex,
% sources/stacks-etale-cohomology/preamble.tex and focused-seeds/AGP/sources/stacks-snapshot-20261011/
% preamble.tex) and those copies are identical in size.  The block below is copied mechanically,
% line for line, from sources/stacks-chow/preamble.tex lines 14-18 and 154-217: the Xy-pic setup and
% the native theorem declarations together with the mathematical macros.  Nothing was added, removed
% or redefined; the wrapper class and the amsmath/amssymb/amsthm/graphicx base packages are supplied
% by the assembler.
% For commutative diagrams we use Xy-pic
% Chapter: stacks-cohomology.tex -- Cohomology of Sheaves.
\usepackage[all]{xy}

% We use 2cell for 2-commutative diagrams.
\xyoption{2cell}

% Theorem environments.
%
\theoremstyle{plain}
\newtheorem{theorem}[subsection]{Theorem}
\newtheorem{proposition}[subsection]{Proposition}
\newtheorem{lemma}[subsection]{Lemma}

\theoremstyle{definition}
\newtheorem{definition}[subsection]{Definition}
\newtheorem{example}[subsection]{Example}
\newtheorem{exercise}[subsection]{Exercise}
\newtheorem{situation}[subsection]{Situation}

\theoremstyle{remark}
\newtheorem{remark}[subsection]{Remark}
\newtheorem{remarks}[subsection]{Remarks}

\numberwithin{equation}{subsection}

% Macros
%
\def\lim{\mathop{\mathrm{lim}}\nolimits}
\def\colim{\mathop{\mathrm{colim}}\nolimits}
\def\Spec{\mathop{\mathrm{Spec}}}
\def\Hom{\mathop{\mathrm{Hom}}\nolimits}
\def\Ext{\mathop{\mathrm{Ext}}\nolimits}
\def\SheafHom{\mathop{\mathcal{H}\!\mathit{om}}\nolimits}
\def\SheafExt{\mathop{\mathcal{E}\!\mathit{xt}}\nolimits}
\def\Sch{\mathit{Sch}}
\def\Mor{\mathop{\mathrm{Mor}}\nolimits}
\def\Ob{\mathop{\mathrm{Ob}}\nolimits}
\def\Sh{\mathop{\mathit{Sh}}\nolimits}
\def\NL{\mathop{N\!L}\nolimits}
\def\CH{\mathop{\mathrm{CH}}\nolimits}
\def\proetale{{pro\text{-}\acute{e}tale}}
\def\etale{{\acute{e}tale}}
\def\QCoh{\mathit{QCoh}}
\def\Ker{\mathop{\mathrm{Ker}}}
\def\Im{\mathop{\mathrm{Im}}}
\def\Coker{\mathop{\mathrm{Coker}}}
\def\Coim{\mathop{\mathrm{Coim}}}

% Boxtimes
%
\DeclareMathSymbol{\boxtimes}{\mathbin}{AMSa}{"02}

%
% Macros for moduli stacks/spaces
%
\def\QCohstack{\mathcal{QC}\!\mathit{oh}}
\def\Cohstack{\mathcal{C}\!\mathit{oh}}
\def\Spacesstack{\mathcal{S}\!\mathit{paces}}
\def\Quotfunctor{\mathrm{Quot}}
\def\Hilbfunctor{\mathrm{Hilb}}
\def\Curvesstack{\mathcal{C}\!\mathit{urves}}
\def\Polarizedstack{\mathcal{P}\!\mathit{olarized}}
\def\Complexesstack{\mathcal{C}\!\mathit{omplexes}}
% \Pic is the operator that assigns to X its picard group, usage \Pic(X)
% \Picardstack_{X/B} denotes the Picard stack of X over B
% \Picardfunctor_{X/B} denotes the Picard functor of X over B
\def\Pic{\mathop{\mathrm{Pic}}\nolimits}
\def\Picardstack{\mathcal{P}\!\mathit{ic}}
\def\Picardfunctor{\mathrm{Pic}}
\def\Deformationcategory{\mathcal{D}\!\mathit{ef}}
\usepackage{cleveref}
\crefname{originallemma}{original lemma}{original lemmas}
\Crefname{originallemma}{Original lemma}{Original lemmas}
\makeatletter
\newlabel{lemma-injective-restriction-surjective}{{01EA}{source}{}{sourcepointer-1}{}}
\newlabel{lemma-injective-restriction-surjective@cref}{{[originallemma][0][]01EA}{[1][0][]source}{}{}{}}
\makeatother
\crefname{originallemma}{original lemma}{original lemmas}
\Crefname{originallemma}{Original lemma}{Original lemmas}
\makeatletter
\newlabel{homology-lemma-long-exact-sequence-cochain}{{0117}{source}{}{sourcepointer-2}{}}
\newlabel{homology-lemma-long-exact-sequence-cochain@cref}{{[originallemma][0][]0117}{[1][0][]source}{}{}{}}
\makeatother
\title{Mayer-Vietoris long exact sequence for a cover by two open subsets}
\author{The Stacks Project Authors}
\date{stacks-project:cohomology.tex:local-snapshot-20261011; GNU FDL 1.2 or later}
\begin{document}
\maketitle
\noindent\textit{Editorial note.} All delivered mathematical text is the author's own, copied from the pinned Stacks Project chapter file cohomology.tex; no statement or proof step was rewritten, reordered or removed. The author's closing sentence omitting functoriality is reproduced as written, as is the bracketed name Mayer-Vietoris in the environment header.\par\medskip
\noindent\textit{Editorial note.} The chapter file pinned-original-source.tex carries only \textbackslash{}input\{preamble\} on its first line, so no preamble text lives inside the chapter. The retained file is a mechanical, line-for-line copy of the project's single shared preamble as staged in the snapshot: sources/stacks-chow/preamble.tex lines 14-18 (Xy-pic setup) and 154-217 (theorem declarations and mathematical macros). Nothing was added, removed or redefined, and the excerpt serves only as preamble, never as a span coordinate.\par\medskip
\noindent\textit{Editorial note.} The same-chapter pointer is located at its \textbackslash{}label line inside the pinned chapter file; the cross-chapter pointer is located in the project's own tag list staged in the snapshot (sources/stacks-chow/tags.txt), at the verbatim line carrying its canonical tag.\par\medskip
\noindent\textit{Allowed standard prerequisites (editorial).}\begin{itemize}
\item Restriction to an open of an injective module is surjective on sections. Same chapter, Lemma labelled lemma-injective-restriction-surjective (source line 715, Stacks tag 01EA); invoked at proof line 786. Supplies surjectivity of (1, -1) on sections of an injective module over U and V onto the intersections, which turns the degreewise sequence of an injective resolution into a short exact sequence of complexes.
\item The long exact cohomology sequence of a short exact sequence of cochain complexes. Homology, Lemma labelled homology-lemma-long-exact-sequence-cochain (canonical Stacks tag 0117); invoked at proof line 799. Supplies the connecting maps and the long exact sequence once the short exact sequence of complexes is in hand.
\end{itemize}\medskip
\ifdefined\hypertarget\hypertarget{sourcepointer-1}{}\fi
\noindent\textit{Source reference.} Original lemma 01EA: sources/stacks-cohomology/pinned-original-source.tex, source lines 715--715. Same-chapter pointer invoked at proof line 786 for surjectivity of (1, -1) on sections of an injective module; the assembly of the short exact sequence of complexes and the conclusion are retained in the counted span.\par\medskip
\ifdefined\hypertarget\hypertarget{sourcepointer-2}{}\fi
\noindent\textit{Source reference.} Original lemma 0117: sources/stacks-chow/tags.txt, source lines 1353--1353. Cross-chapter standard prerequisite invoked once at proof line 799 to take the long exact cohomology sequence; the short exact sequence of complexes it is applied to is retained in the counted span.\par\medskip
\begin{lemma}[Mayer-Vietoris]
\label{lemma-mayer-vietoris}
Let $X$ be a ringed space. Suppose that $X = U \cup V$ is a
union of two open subsets. For every $\mathcal{O}_X$-module $\mathcal{F}$
there exists a long exact cohomology sequence
$$
0 \to
H^0(X, \mathcal{F}) \to
H^0(U, \mathcal{F}) \oplus H^0(V, \mathcal{F}) \to
H^0(U \cap V, \mathcal{F}) \to
H^1(X, \mathcal{F}) \to \ldots
$$
This long exact sequence is functorial in $\mathcal{F}$.
\end{lemma}

\begin{proof}
The sheaf condition says that the kernel of
$(1, -1) : \mathcal{F}(U) \oplus \mathcal{F}(V) \to \mathcal{F}(U \cap V)$
is equal to the image of $\mathcal{F}(X)$ by the first map
for any abelian sheaf $\mathcal{F}$.
Lemma \ref{lemma-injective-restriction-surjective} above implies that the map
$(1, -1) : \mathcal{I}(U) \oplus \mathcal{I}(V) \to \mathcal{I}(U \cap V)$
is surjective whenever $\mathcal{I}$ is an injective $\mathcal{O}_X$-module.
Hence if $\mathcal{F} \to \mathcal{I}^\bullet$ is an injective resolution
of $\mathcal{F}$, then we get a short exact sequence of complexes
$$
0 \to
\mathcal{I}^\bullet(X) \to
\mathcal{I}^\bullet(U) \oplus \mathcal{I}^\bullet(V) \to
\mathcal{I}^\bullet(U \cap V) \to
0.
$$
Taking cohomology gives the result (use
Homology, Lemma \ref{homology-lemma-long-exact-sequence-cochain}).
We omit the proof of the functoriality of the sequence.
\end{proof}

\end{document}
